% Part: incompleteness
% Chapter: incompleteness-provability
% Section: introduction

\documentclass[../../../include/open-logic-section]{subfiles}

\begin{document}

\olfileid{inc}{inp}{int}

\olsection{પરિચય}

હિલ્બર્ટ માનતા હતા કે
પ્રાકૃતિક સંખ્યાઓ જેવી
ગણિતીય રચના માટેની
સ્વયંસિદ્ધિ-પદ્ધતિ તે
રચના વિશેનાં બધાં સાચાં
વિધાનો તારવી ન શકે તો
અપૂરતી છે. પછી તેમને
ઔપચારિક નિગમન-પદ્ધતિઓમાં
રસ પડ્યો, તેની સાથે આ
વિચાર જોડીએ તો એમ લાગે
છે કે પ્રાકૃતિક સંખ્યાઓ
વિશે તર્ક કરવા વપરાતી
ઔપચારિક પદ્ધતિ માત્ર
સુસંગત જ નહીં, પણ
\emph{પૂર્ણ} પણ હોવી જોઈએ
એવું તેઓ માનતા હતા:
એટલે તેની ભાષામાંનું દરેક
વિધાન અથવા તેનો નિષેધ
!!{derivable} હોય.
ગોડેલનું પ્રથમ અપૂર્ણતા
પ્રમેય બતાવે છે કે આવી
સ્વયંસિદ્ધિ-પદ્ધતિ નથી:
અંકગણિત માટે પૂર્ણ,
સુસંગત અને
!!{axiomatizable}
ઔપચારિક પદ્ધતિ નથી.
હકીકતમાં ``પૂરતા મજબૂત'',
સુસંગત અને
!!{axiomatizable}
ગણિતીય સિદ્ધાંતોમાંથી
કોઈપણ પૂર્ણ નથી.

હિલ્બર્ટનું એક વધુ મહત્વનું
ધ્યેય, આધુનિક
(``શાસ્ત્રીય'') ગણિતને
વાજબી ઠેરવવાના તેમના
કાર્યક્રમનું કેન્દ્ર, શાસ્ત્રીય
તર્કને રજૂ કરતી ઔપચારિક
પદ્ધતિઓ માટે સાન્તવાદી
સુસંગતતા-સાબિતીઓ શોધવાનું
હતું. તેથી હિલ્બર્ટના
કાર્યક્રમ માટે ગોડેલનું
બીજું અપૂર્ણતા પ્રમેય
ઘણો મોટો આઘાત હતો.
બીજા અપૂર્ણતા પ્રમેયને પણ
પહેલાની જેમ અચોક્કસ
શબ્દોમાં કહી શકાય.
આશરે તે કહે છે કે
અંકગણિતનો કોઈ પૂરતો
મજબૂત સિદ્ધાંત પોતાની
સુસંગતતા સાબિત કરી
શકતો નથી. ``પૂરતો
મજબૂત''માં $\Th{Q}$થી
થોડું વધુ આવરી લેવું
પડશે.

અપૂર્ણતા પ્રમેયની ગોડેલની
મૂળ સાબિતીનો વિચાર
એપિમેનિડીસના વિરોધાભાસમાં
દેખાય છે. ક્રિટનો વતની
એપિમેનિડીસે કહ્યું કે
ક્રિટના બધા લોકો જૂઠા છે;
વિરોધાભાસનું વધુ સીધું
રૂપ ``આ વાક્ય અસત્ય છે''
એવું વિધાન છે. સત્યને
!!{derivability}થી બદલીને
ગોડેલ એવા !!a{sentence}ને
ઔપચારિક બનાવી શક્યા જે
આડકતરી રીતે પોતે
!!{derivable} નથી એમ
કહે છે. જો તે !!{sentence}
!!{derivable} હોય, તો
સિદ્ધાંત અસુસંગત બને.
ગોડેલે બતાવ્યું કે તેઓ
વિચારતા હતા એવી
સ્વયંસિદ્ધિ-પદ્ધતિમાંથી
તે !!{sentence}નો નિષેધ
પણ !!{derivable} નથી.
(આ બીજા ભાગ માટે ગોડેલે
સિદ્ધાંત~$\Th{T}$ને
``$\omega$-સુસંગત''
હોવાની ધારણા કરવી પડી.
$\omega$-સુસંગતતા
સુસંગતતા સાથે સંબંધિત છે,
પરંતુ વધુ મજબૂત ગુણધર્મ
છે.\footnote{એટલે દરેક
$\omega$-સુસંગત સિદ્ધાંત
સુસંગત હોય છે, પરંતુ
દરેક સુસંગત સિદ્ધાંત
એવો હોય જ એવું
નથી.} ગોડેલના થોડાં વર્ષો
પછી રોસરે બતાવ્યું કે
માત્ર~$\Th{T}$ની
સુસંગતતા ધારવી પૂરતી છે.)

પહેલો પડકાર એવો
!!a{sentence} કેવી રીતે
રચવું તે સમજવાનો છે,
જે પોતાનો જ ઉલ્લેખ કરે.
$\Th{Q}$ની ભાષાના દરેક
!!{formula}~$!A$ માટે
~$\gn{!A}$ એ
~$\Gn{!A}$ને અનુરૂપ
અંકપદ દર્શાવે એમ માનો.
આનો અર્થ વિચારો:
$!A$ એ~$\Th{Q}$ની
ભાષાનું !!a{formula} છે,
$\Gn{!A}$ એ પ્રાકૃતિક
સંખ્યા છે, અને
$\gn{!A}$ એ~$\Th{Q}$ની
ભાષાનું \emph{પદ} છે.
આથી $\Th{Q}$ની ભાષાના
દરેક !!{formula}~$!A$ને
એક \emph{નામ},
~$\gn{!A}$, મળે છે,
જે $\Th{Q}$ની ભાષાનું
પદ છે. આમ એક વૈચારિક
માળખું મળે છે જેમાં
$\Th{Q}$ની ભાષાનાં
!!{formula} બીજા
!!{formula} વિશે કંઈક
``કહી'' શકે છે.
નીચેનું લેમ્મા નિશ્ચિત-બિંદુ
લેમ્મા કહેવાય છે.

\begin{lem}
$\Th{T}$ એ~$\Th{Q}$નું
કોઈપણ વિસ્તરણ હોય અને
$!B(x)$ એવું કોઈપણ
!!{formula} હોય જેમાં
માત્ર ચલ~$x$ મુક્ત હોય.
તો એવું !!a{sentence}~$!A$
છે કે
$\Th{T} \Proves!A \liff !B(\gn{!A})$.
\end{lem}

લેમ્મા કહે છે કે કોઈપણ
ગુણધર્મ~$!B(x)$ આપેલો
હોય, તો એવું
!!a{sentence}~$!A$ છે
જે ``$!B(x)$ મારા માટે
સાચું છે'' એમ કહે છે,
અને $\Th{T}$ પણ આ
``જાણે'' છે.

આવું !!a{sentence} કેવી
રીતે રચી શકીએ?
ક્વાઇનનું નીચેનું
એપિમેનિડીસ-વિરોધાભાસનું
રૂપ વિચારો:
\begin{quote}
``પોતાનું અવતરણ પહેલાં
મૂકવાથી અસત્ય મળે છે''
પોતાનું અવતરણ પહેલાં
મૂકવાથી અસત્ય મળે છે.
\end{quote}
આ વાક્ય સીધો પોતાનો
ઉલ્લેખ કરતું નથી.
તે અવતરણચિહ્નો વચ્ચેના
વાક્યખંડ નામના વાક્યરચનાત્મક
પદાર્થ વિશે માત્ર વિધાન
કરે છે; આમ તે નીચેનાં
વાક્યોની હરોળમાં આવે છે:
\begin{enumerate}
\item ``રોબર્ટ'' સુંદર નામ છે.
\item ``હું દોડ્યો.'' ટૂંકું વાક્ય છે.
\item ``ત્રણ શબ્દો છે''માં ત્રણ શબ્દો છે.
\end{enumerate}
પરંતુ ``પોતાનું અવતરણ
પહેલાં મૂકવાથી અસત્ય મળે
છે'' વાક્યખંડની પહેલાં
એ જ વાક્યખંડને
અવતરણચિહ્નોમાં મૂકીશું
તો શું થશે? આપણને
મૂળ વાક્ય જ મળશે!{}
ટૂંકમાં, આ વાક્ય પોતે
અસત્ય છે એમ કહે છે.

\end{document}
